% TEMPLATE-MILC-v3.tex - plantilla maestra UNICA de las guias MILC v3.
%
% La usan las tres familias de guias, cada una con su driver:
%   · Grados 6-11  -> scripts/build-guias-g11.py
%   · Bebras       -> scripts/build-guias-bebras.py
%   · Semillero    -> scripts/build-guias-semillero.py
%
% Antes habia tres copias casi identicas (~400 lineas iguales, 6 distintas):
% cada mejora habia que aplicarla tres veces, y en la practica no se hacia
% (el motor de recursos vivia solo en dos de ellas). Lo que varia entre
% familias esta parametrizado en 6 placeholders que cada driver llena
% (se nombran aqui SIN su sintaxis real: el builder reemplaza por texto plano,
% tambien dentro de los comentarios, y un valor multilinea rompe el preambulo):
%
%   HEADER_IZQ        encabezado izquierdo de pagina
%   HEADER_DER        encabezado derecho de pagina
%   PORTADA_ETIQUETA  rotulo sobre el numero grande (GRADO/MOMENTO/MODULO)
%   PORTADA_SERIE     linea de serie de la portada
%   PORTADA_ID        identificador de la guia en la portada
%   PORTADA_CIRCULO   el circulito de grado (°), o vacio si no aplica
%
% El resto de claves las documenta content/guias/_SCHEMA.md. El script valida
% que no queden placeholders sin reemplazar antes de escribir el archivo final.

\documentclass[11pt,letterpaper]{article}
\usepackage[margin=2.0cm]{geometry}
\usepackage[table]{xcolor}
\usepackage{fontspec}
\usepackage[spanish]{babel}
\usepackage{tikz}
% Librerías TikZ para los diagramas de las guías (bloque `recursos:`):
%  · babel  → babel en español vuelve activos caracteres como «>», y TikZ los usa
%             en sus opciones (>=stealth, ->). Sin esta librería, cualquier
%             diagrama con flechas revienta con «\language@active@arg> has an
%             extra }». Debe ir ANTES de que se dibuje nada.
%  · positioning        → sintaxis «below left=2pt and 3pt».
%  · decorations.pathreplacing → llaves ({brace}) para señalar rangos.
\usetikzlibrary{babel,positioning,decorations.pathreplacing}
\usepackage{graphicx}
% Circuitos: el trabajo de electrónica se hace hoy con micro:bit y sus sensores
% integrados (MakeCode, sin cableado), así que no se necesitan esquemáticos.
% Si algún día entran montajes con protoboard, basta descomentar esta línea:
% los diagramas .tex/.tikz declarados en `recursos:` ya se incrustan solos.
% \usepackage{circuitikz}
\usepackage[most]{tcolorbox}
\usepackage{tabularx}
\usepackage{array}
\usepackage{fancyhdr}
\usepackage{titlesec}
\usepackage[document]{ragged2e}  % bandera a la derecha en todo el cuerpo
\usepackage{enumitem}
\usepackage{microtype}

% Helvetica Neue no trae la flecha U+2192, asi que cada «→» escrito literalmente
% en el YAML se compilaba a NADA, en silencio (462 flechas en 61 guias: rutas del
% tipo «paleta → tipografia → lockup» salian rotas). Mapeamos el caracter a su
% version matematica, que si tiene glifo. Asi el autor puede seguir escribiendo
% «→» comodamente en el YAML.
\usepackage{newunicodechar}
\newunicodechar{→}{\ensuremath{\rightarrow}}
% Mismo problema con los simbolos de marca y vinetas: el «✓» de «una palomita ✓»
% salia como cuadrito vacio en el PDF de todas las guias que lo usaban.
\usepackage{pifont}
\newunicodechar{✓}{\ding{51}}
\newunicodechar{✗}{\ding{55}}
\newunicodechar{✔}{\ding{52}}
\newunicodechar{•}{\textbullet}
\newunicodechar{≥}{\ensuremath{\geq}}
\newunicodechar{≤}{\ensuremath{\leq}}

\setmainfont{Helvetica Neue}
\setsansfont{Helvetica Neue}
\setlength{\parindent}{0pt}
\setlength{\parskip}{6pt}
% Sin particion de palabras: en 6.º-7.º una palabra cortada por guion al final
% de linea («Comuni-cacion») frena la lectura mucho mas que una linea corta.
\hyphenpenalty=10000
\exhyphenpenalty=10000
\pretolerance=2000
\tolerance=3000
\setlength{\emergencystretch}{3em}
\linespread{1.10}
\renewcommand{\arraystretch}{1.25}
\setlist{leftmargin=5mm,itemsep=1mm,topsep=1mm}
\newcolumntype{Y}{>{\RaggedRight\arraybackslash}X}

\definecolor{milcMagenta}{HTML}{E6007E}
\definecolor{milcVino}{HTML}{5A0038}
\definecolor{milcTurquesa}{HTML}{0093A5}
\definecolor{milcVerde}{HTML}{58A923}
\definecolor{milcMostaza}{HTML}{E5B400}
\definecolor{milcOcre}{HTML}{B86600}
\definecolor{milcNegro}{HTML}{241816}
\definecolor{milcGris}{HTML}{F7F7F4}
\definecolor{milcLinea}{HTML}{E8E2DD}
\definecolor{milcGradePrimary}{HTML}{0066FF}
\definecolor{milcGradeDark}{HTML}{003D99}
\definecolor{milcGradeSoft}{HTML}{D6E8FF}
\definecolor{milcGradeAccent}{HTML}{CFE7FF}
\definecolor{milcGradeText}{HTML}{FFFFFF}
\color{milcNegro}

\pagestyle{fancy}
\fancyhf{}
\lhead{\small Tecnología e Informática · Grado 11}
\rhead{\small Guía 13 · MILC v3}
\cfoot{\small\thepage}
\renewcommand{\headrulewidth}{0pt}

\newtcolorbox{titlebox}[1]{
  enhanced,colback=#1,colframe=#1,arc=7mm,boxrule=0pt,
  left=7mm,right=7mm,top=3mm,bottom=3mm,
  before skip=8pt,after skip=8pt,
  fontupper=\sffamily\bfseries\Large\color{white}
}

\newtcolorbox{softbox}[3]{
  enhanced,colback=#2,colframe=#1,borderline west={4pt}{0pt}{#1},
  arc=2mm,boxrule=.4pt,left=4mm,right=4mm,top=3mm,bottom=3mm,
  before skip=6pt,after skip=8pt,title={#3},coltitle=white,
  colbacktitle=#1,fonttitle=\sffamily\bfseries,fontupper=\RaggedRight,
  attach boxed title to top left={xshift=3mm,yshift=-2mm},
  boxed title style={arc=2mm, boxrule=0pt}
}

\newtcolorbox{drawbox}[2]{
  enhanced,colback=white,colframe=#1,arc=4mm,boxrule=1pt,
  left=5mm,right=5mm,top=4mm,bottom=4mm,before skip=8pt,after skip=8pt,
  title={#2},coltitle=white,colbacktitle=#1,fonttitle=\sffamily\bfseries,
  fontupper=\RaggedRight,
  attach boxed title to top left={xshift=4mm,yshift=-2mm},
  boxed title style={arc=3mm, boxrule=0pt}
}

\newtcolorbox{infoband}[1]{
  enhanced,colback=#1!8,colframe=#1!50,arc=2mm,boxrule=.5pt,
  left=4mm,right=4mm,top=2mm,bottom=2mm,before skip=4pt,after skip=4pt,
  fontupper=\small\RaggedRight
}

% Puente narrativo entre secciones (contrato editorial Conéctate).
% Da continuidad: "Acabas de X. Ahora vas a Y porque Z."
% Visualmente: banda izquierda en color del grado, texto en itálica pequeña.
\newtcolorbox{puentebox}{
  enhanced,colback=milcGradeSoft,colframe=milcGradeDark,
  borderline west={3pt}{0pt}{milcGradeDark},
  arc=0mm,boxrule=0pt,left=5mm,right=4mm,top=2.5mm,bottom=2.5mm,
  before skip=4pt,after skip=8pt,
  fontupper=\itshape\small\color{milcNegro}\RaggedRight
}
% La flecha va en modo matematico a proposito: Helvetica Neue NO tiene el glifo
% U+2192, asi que \textrightarrow se compilaba a NADA («Missing character: There
% is no → in font Helvetica Neue») y el puente salia sin flecha. En modo
% matematico la flecha viene de la fuente matematica y si se dibuja.
\newcommand{\puente}[1]{%
  \begin{puentebox}%
    \textcolor{milcGradeDark}{$\rightarrow$}\hspace{2mm}#1%
  \end{puentebox}%
}

\newcommand{\linead}[1]{\noindent\color{milcLinea}\rule{#1}{0.5pt}\color{milcNegro}}
\newcommand{\respuesta}[1]{\par\vspace{2mm}\foreach \n in {1,...,#1}{\noindent\color{milcLinea}\rule{\linewidth}{0.5pt}\par\vspace{5mm}}\color{milcNegro}}
\newcommand{\checkbox}{\textcolor{milcGradeDark}{\fbox{\phantom{x}}}\hspace{5pt}}

\newcommand{\phasecard}[4]{
\begin{tcolorbox}[enhanced,colback=#3,colframe=#2,arc=3mm,boxrule=.5pt,height=3.05cm,valign=center,halign=center,left=1.5mm,right=1.5mm,top=2mm,bottom=2mm]
{\sffamily\bfseries\fontsize{22}{24}\selectfont\textcolor{#2}{#1}}\par
{\sffamily\bfseries\small #4}
\end{tcolorbox}
}

% ── Piezas del rediseno 2026 (legibilidad 6.º-11.º) ───────────────────
% Banda de estacion: reemplaza el titlebox macizo. Titulo a la izquierda,
% dato practico (tiempo/modalidad) a la derecha, en una sola linea.
\newcommand{\milcestacion}[2]{%
\begin{tcolorbox}[enhanced,colback=#1,colframe=#1,arc=2mm,boxrule=0pt,
  left=5mm,right=5mm,top=2.6mm,bottom=2.6mm,before skip=11pt,after skip=7pt]
{\sffamily\bfseries\fontsize{15}{18}\selectfont\color{white}#2}
\end{tcolorbox}}

% Tarjeta de paso numerada: sustituye las tablas «Pilar / Aplicacion», que
% eran formato de consulta docente, no de estudiante. El numero va en un
% circulo sobre el borde para que la secuencia se lea de un vistazo.
\newcommand{\milcpaso}[3]{%
\begin{tcolorbox}[enhanced,colback=white,colframe=milcLinea,arc=1.5mm,boxrule=.9pt,
  left=14mm,right=4mm,top=3mm,bottom=3mm,before skip=4pt,after skip=4pt,
  fontupper=\RaggedRight,
  overlay={\node[anchor=north west,circle,fill=#2,text=white,inner sep=0pt,
    minimum size=7.4mm,font=\sffamily\bfseries\small]
    at ([xshift=3.6mm,yshift=-3.2mm]frame.north west) {#1};}]
#3
\end{tcolorbox}}

% Casilla marcable de tamano util con lapiz (la anterior era \fbox{\phantom{x}},
% de alto variable segun el texto que la seguia).
\newcommand{\milcmarca}{\framebox[3.8mm]{\rule{0pt}{3.4mm}}}

% Fila marcable de lista suelta (checks de la estacion 1).
\newcommand{\milccheckrow}[1]{%
\par\vspace{1.8mm}\noindent
\makebox[6.4mm][l]{\raisebox{-.55mm}{\textcolor{milcNegro}{\milcmarca}}}%
\parbox[t]{\dimexpr\linewidth-6.4mm\relax}{\RaggedRight #1}\par}

% Celda marcable dentro de la rubrica (tabla de autoevaluacion). Misma
% sangria francesa, pero pensada para vivir dentro de una columna X.
\newcommand{\milcopcion}[1]{%
\makebox[5.2mm][l]{\raisebox{-.55mm}{\textcolor{milcNegro}{\milcmarca}}}%
\parbox[t]{\dimexpr\linewidth-5.2mm\relax}{\RaggedRight #1}}

% ── Recursos visuales de la guía ──────────────────────────────────────
% Declarados en el bloque `recursos:` del YAML y emitidos por el builder.
% \guiaFigura   → imagen rasterizada o PDF (foto, carta celeste, montaje).
% \guiaDiagrama → fuente TikZ/circuitikz (.tex) incrustada como vector:
%                 es la vía para circuitos y esquemáticos editables.
% #1 = ruta relativa al .tex · #2 = pie (puede ir vacío)
\newcommand{\guiaFigura}[2]{%
\begin{center}
\includegraphics[width=0.88\linewidth,height=0.38\textheight,keepaspectratio]{#1}\\[1.5mm]
{\footnotesize\itshape\textcolor{milcNegro!75}{#2}}
\end{center}
}

\newcommand{\guiaDiagrama}[2]{%
\begin{center}
\input{#1}\\[1.5mm]
{\footnotesize\itshape\textcolor{milcNegro!75}{#2}}
\end{center}
}

\begin{document}
\thispagestyle{empty}

% ── Portada ───────────────────────────────────────────────────────────
% Antes: un rectangulo de color a pagina completa. Se veia bien en pantalla,
% pero en la fotocopiadora del colegio se comia el toner y salia gris sucio.
% Ahora la banda ocupa el tercio superior y el resto va en blanco.
\begin{tikzpicture}[remember picture,overlay]
  \path[fill=milcGradePrimary]
    (current page.north west) rectangle ([yshift=-10.6cm]current page.north east);

  \node[anchor=north west,text=milcGradeText] at ([xshift=2.0cm,yshift=-1.55cm]current page.north west)
    {\sffamily\bfseries\fontsize{12}{14}\selectfont GRADO};
  \node[anchor=north west,text=milcGradeText,opacity=.85] at ([xshift=2.0cm,yshift=-2.35cm]current page.north west)
    {\sffamily\bfseries\fontsize{8.5}{10}\selectfont SERIE GUÍAS · TECNOLOGÍA E INFORMÁTICA};
  \node[anchor=north east,text=milcGradeText,opacity=.85,align=right] at ([xshift=-2.0cm,yshift=-1.55cm]current page.north east)
    {\sffamily\bfseries\fontsize{9}{12}\selectfont I.E. Sor María Juliana\\Cartago, Valle del Cauca};

  \node[anchor=west,text=milcGradeText] (gradonum) at ([xshift=1.85cm,yshift=-7.1cm]current page.north west)
    {\sffamily\bfseries\fontsize{112}{112}\selectfont 11};
  % El circulito de grado (°) solo aplica a las guias de grado («11°»); en
  % Bebras y Semillero el numero grande es un momento/modulo, no un grado.
  % Se posiciona relativo al nodo `gradonum`, no en coordenadas absolutas.
  \draw[milcGradeText,line width=1.6pt]
    ([xshift=2.0mm,yshift=-4.5mm]gradonum.north east) circle (.15cm);

  \node[anchor=west,text=milcGradeText,text width=12.3cm,align=left] at ([xshift=6.5cm,yshift=-6.6cm]current page.north west)
    {
     {\sffamily\bfseries\fontsize{9}{11}\selectfont\MakeUppercase{Período 2 · Automatización y procesos}}\\[3mm]
     {\sffamily\bfseries\fontsize{21}{25}\selectfont Formularios digitales ---\\captura sin\\papel}\\[3.5mm]
     {\sffamily\bfseries\fontsize{15}{18}\selectfont Guía 13-11-TIC}
    };
\end{tikzpicture}

\vspace*{9.4cm}
\vfill

{\sffamily\bfseries\fontsize{8}{10}\selectfont\textcolor{milcGradeDark}{LO QUE TE LLEVAS HOY}}

\begin{tcolorbox}[enhanced,colback=white,colframe=milcNegro,arc=2mm,boxrule=1.6pt,
  left=5mm,right=5mm,top=4mm,bottom=4mm,before skip=3pt,after skip=12pt,
  fontupper=\RaggedRight\fontsize{11}{15.5}\selectfont]
1 formulario funcional (Google Forms o Typeform) con al menos 8 campos, validaciones activas (obligatorio, formato email, rango numérico), lógica condicional mínima (1 ramificación) y respuestas que llegan a una hoja de cálculo enlazada, diseñado para un proceso real
\end{tcolorbox}

\begin{tcolorbox}[enhanced,colback=milcGris,colframe=milcGris,arc=2mm,boxrule=0pt,
  left=5mm,right=5mm,top=3.5mm,bottom=3.5mm,after skip=12pt]
{\sffamily\bfseries\fontsize{8}{10}\selectfont\textcolor{milcNegro!55}{ESTA GUÍA ES DE}}\\[3mm]
\begin{tabularx}{\linewidth}{X p{3.2cm} p{3.2cm}}
\linead{\linewidth} & \linead{3.0cm} & \linead{3.0cm}\\[-1mm]
{\fontsize{8.5}{10}\selectfont\textcolor{milcNegro!55}{Nombre completo}} &
{\fontsize{8.5}{10}\selectfont\textcolor{milcNegro!55}{Grupo}} &
{\fontsize{8.5}{10}\selectfont\textcolor{milcNegro!55}{Fecha}}\\
\end{tabularx}
\end{tcolorbox}

\vfill

\noindent\textcolor{milcLinea}{\rule{\linewidth}{1pt}}\\[2mm]
{\sffamily\bfseries\fontsize{9.5}{12}\selectfont Autor: Dr. Álvaro Cárdenas Orozco}\\[1mm]
{\sffamily\fontsize{8.5}{11}\selectfont\textcolor{milcNegro!55}{Metodología MILC · apertura ancestral · pensamiento computacional · triángulo Dussel–estoicismo–Floridi}}

\newpage

\section*{Apertura: Formularios digitales --- captura de información sin papel}

\begin{softbox}{milcGradeDark}{milcGradeSoft}{Saber ancestral}
La \textbf{libreta del jornal} de las fincas cafeteras del Valle es un formulario de papel: la abuela apunta cuántos canastos cosechó cada jornalero, cuántas horas trabajó, cuánto se pagó. La \textbf{libreta del fiado} del tendero del barrio registra quién pidió qué y cuándo. El \textbf{cuaderno de matrícula} del docente apunta nombre, acudiente, dirección, fecha. Estos cuadernos comparten una sabiduría: \emph{capturar información estructurada, una sola vez, en el momento exacto en que ocurre}. Los formularios digitales son la libreta moderna --- pero con cálculo automático, sin tachones, sin pérdida de hojas y con respaldo en la nube. El gesto antiguo de \emph{``apuntar mientras todavía está fresco''} sigue siendo el mismo.
\end{softbox}

\begin{softbox}{milcTurquesa}{milcGris}{Saber contemporáneo}
Un \textbf{formulario digital} es una herramienta para capturar información estructurada de forma uniforme. Sus elementos básicos son: (1) \textbf{campos} (texto corto, texto largo, opción múltiple, casillas, escala, fecha, email, número); (2) \textbf{validaciones} (obligatorio, formato esperado, rango numérico); (3) \textbf{lógica condicional} (\emph{``si responde A, muestra la pregunta 5; si responde B, salta a la 8''}); (4) \textbf{destino} (hoja de cálculo enlazada que almacena las respuestas). Las dos herramientas más comunes en Colombia son \texttt{Google Forms} (gratuito, integrado con Drive) y \texttt{Typeform} (más visual, plan gratuito limitado). Diseñar un formulario es \emph{diseño de experiencia}: cada pregunta innecesaria hace que el usuario abandone.
\end{softbox}

\begin{softbox}{milcMostaza}{milcGris}{Pregunta-puente}
\textit{¿Cómo sabía la abuela qué campos eran imprescindibles en su libreta del jornal y cuáles sobraban? ¿Y qué pierde un emprendedor novato cuando inventa un formulario de 25 preguntas sin haber mapeado primero el proceso que pretende digitalizar?}
\end{softbox}

\begin{softbox}{milcVerde}{milcGris}{Saber y saber hacer}
Al terminar podrás: (1) \textbf{identificar} los 8 tipos básicos de campo de un formulario digital y cuándo usar cada uno; (2) \textbf{analizar} qué información necesita realmente un proceso para no pedir datos innecesarios; (3) \textbf{crear} un formulario funcional con validaciones y lógica condicional sobre un proceso real; (4) \textbf{evaluar} la usabilidad del formulario respondiéndolo tú mismo y midiendo cuánto demoró.
\end{softbox}



\small
\begin{tabularx}{\textwidth}{>{\bfseries}p{3.0cm}Y}
\rowcolor{milcGradePrimary!12}Autor & Dr. Álvaro Cárdenas Orozco\\
DBA & Diseño y mejora de procesos digitales con criterio ético (MEN, T\&I)\\
Referentes & Pensamiento computacional · Lean / Kaizen · Floridi (ética informacional)\\
Producto final & 1 formulario funcional (Google Forms o Typeform) con al menos 8 campos, validaciones activas (obligatorio, formato email, rango numérico), lógica condicional mínima (1 ramificación) y respuestas que llegan a una hoja de cálculo enlazada, diseñado para un proceso real\\
\end{tabularx}
\normalsize

\puente{Antes de empezar, mira el viaje completo. En la sesión 1 mapeaste un proceso real; en la sesión 2 lo diagramaste en BPMN. Hoy automatizas el \emph{primer paso} de ese proceso: la captura de información. Es el inicio práctico de la automatización.}

\milcestacion{milcGradeDark}{Ruta MILC: así vamos a aprender}
\begin{tabularx}{\textwidth}{>{\centering\arraybackslash}X>{\centering\arraybackslash}X>{\centering\arraybackslash}X>{\centering\arraybackslash}X}
\phasecard{1}{milcVerde}{milcGris}{Escucha\\[-1mm]\small Observo, escucho y pregunto.} &
\phasecard{2}{milcTurquesa}{milcGris}{Entender\\[-1mm]\small Sistematizo y ordeno.} &
\phasecard{3}{milcMagenta}{milcGris}{Hacer\\[-1mm]\small Construyo y aplico.} &
\phasecard{4}{milcVino}{milcGris}{Cierre\\[-1mm]\small Evalúo y reflexiono.}
\end{tabularx}


\puente{Antes de teorizar, vas a responder 2 formularios reales (uno bien hecho, uno mal hecho) y a anotar qué te frustró y qué te facilitó. Esa experiencia como usuario te enseña qué evitar como diseñador.}

\milcestacion{milcVerde}{Estación 1 de 3 · Escucha}

\begin{softbox}{milcVerde}{milcGris}{Escena para escuchar}
\textbf{Actividad 1 · IDENTIFICA --- Responde 2 formularios reales} (15 min · individual). El docente comparte enlaces a 2 formularios reales: uno bien diseñado (corto, claro, con campos pertinentes), otro mal diseñado (largo, con preguntas redundantes, sin validaciones). Respóndelos ambos y anota: cuánto demoraste, qué te frustró, qué te facilitó, en cuál abandonaste mentalmente y en cuál no. Esa experiencia como usuario es la mejor escuela de diseño.
\end{softbox}

{\sffamily\bfseries\fontsize{8}{10}\selectfont\textcolor{milcNegro!55}{MARCA CUANDO LO HAYAS HECHO}}
\milccheckrow{Respondí los 2 formularios completos y medí tiempo.}
\milccheckrow{Anoté qué me frustró y qué me facilitó en cada uno.}
\milccheckrow{Identifiqué al menos 2 diferencias concretas de diseño entre el bueno y el malo.}

\begin{infoband}{milcVerde}


\textbf{Lo que va al cuaderno (Actividad 1 · IDENTIFICA).} Título: «Actividad 1 --- Responde 2 formularios». \textbf{Formato:} tabla 2 filas (un formulario por fila) × 3 columnas (Tiempo~|~Lo que me frustró~|~Lo que me facilitó). Al final 2 líneas con las diferencias clave de diseño. \textbf{Sabes que terminaste cuando} podrías describir a un amigo qué hace bueno o malo a un formulario sin recurrir a vocabulario técnico.
\end{infoband}

\puente{Ya viste el contraste. Ahora vas a entender la \textbf{anatomía} de un formulario bien diseñado: qué tipos de campo elegir, qué validaciones poner, cuándo usar lógica condicional, cómo evitar la fatiga del usuario, cómo enlazar a hoja de cálculo.}

\milcestacion{milcTurquesa}{Estación 2 de 3 · Entender}

\begin{softbox}{milcTurquesa}{milcGris}{Lo que vas a entender}
Diseñar un formulario parece trivial pero tiene reglas de oficio. Vamos a descomponerlas con los pilares del pensamiento computacional para que cada decisión de diseño tenga justificación, no preferencia personal.
\end{softbox}

\milcpaso{1}{milcTurquesa}{El formulario se descompone en 4 capas: (1) \textbf{Campos}: los 8 tipos básicos son \emph{texto corto} (nombre, ciudad), \emph{texto largo} (comentario), \emph{opción múltiple} (una entre varias), \emph{casillas} (varias entre varias), \emph{escala} (1-5), \emph{fecha}, \emph{email}, \emph{número}. (2) \textbf{Validaciones}: obligatorio (no avanza sin respuesta), formato esperado (email con \@), rango numérico (entre 0 y 100). (3) \textbf{Lógica condicional}: \emph{si responde X, salta a la pregunta Y}. (4) \textbf{Destino}: la hoja de cálculo enlazada que recibe las respuestas.}
\milcpaso{2}{milcTurquesa}{Todo formulario sigue el patrón \textbf{``cada pregunta paga renta''}: si una pregunta no es indispensable para el proceso, se elimina. Cada pregunta innecesaria aumenta la tasa de abandono. Regla operativa: \emph{si la respuesta no se va a usar para tomar una decisión o ejecutar una acción, no la pidas}. El diseñador profesional cuida los minutos del usuario como su tiempo propio.}
\milcpaso{3}{milcTurquesa}{Lo esencial cabe en una pregunta: \emph{¿en qué hoja de cálculo van a aparecer estas respuestas y quién las va a procesar?}. Si la respuesta es \emph{``no sé''}, el formulario está mal diseñado. Todo formulario debe tener destino claro: hoja, persona responsable, frecuencia de revisión. Sin destino, las respuestas se acumulan sin acción.}
\milcpaso{4}{milcTurquesa}{Algoritmo: (1) define el proceso que automatizas; (2) lista qué información necesitas para ejecutarlo; (3) elige el tipo de campo apropiado para cada dato; (4) decide qué campos son obligatorios; (5) agrega lógica condicional si el flujo tiene ramificaciones; (6) enlaza a hoja de cálculo; (7) responde tú mismo el formulario; (8) ajusta lo que te trabó como usuario.}

\begin{drawbox}{milcTurquesa}{Anatomía de un formulario bien diseñado}
\textbf{Título:} qué proceso atiende. \\[1mm] \textbf{Descripción:} 1 frase con para qué se usan las respuestas. \\[1mm] \textbf{Campos:} 5-10 preguntas máximo (más, abandono garantizado). \\[1mm] \textbf{Validaciones:} obligatorio + formato esperado en campos críticos. \\[1mm] \textbf{Lógica:} ramificaciones si el flujo lo exige. \\[1mm] \textbf{Destino:} hoja de cálculo enlazada + responsable de revisión. \\[1mm] \textbf{Confirmación:} mensaje claro al usuario al enviar.
\end{drawbox}

\begin{softbox}{milcMostaza}{milcGris}{Errores comunes y su corrección}
(1) Hacer un formulario de 25 preguntas \emph{``por si acaso''}: 80\% de los usuarios abandona después de la pregunta 10. (2) Olvidar validaciones: la columna \emph{email} llega con \emph{``juancarlos''} en vez de \emph{``juan@correo.com''}. (3) No enlazar a hoja: las respuestas se quedan dentro del formulario y nadie las usa. (4) Preguntas redundantes: pedir nombre completo, después nombres, después apellidos por separado. (5) Diseñar sin probar como usuario: el creador no detecta los problemas que cualquier usuario percibe en 30 segundos.
\end{softbox}

\begin{infoband}{milcTurquesa}


\textbf{Lo que va al cuaderno (Actividad 2 · ANALIZA).} Título: «Actividad 2 --- Anatomía del formulario». \textbf{Formato:} ficha con las 7 partes del formulario aplicadas al proceso que vas a automatizar + 1 lista de los 5 errores comunes con marca del que más temes cometer. \textbf{Sabes que terminaste cuando} podrías auditar un formulario ajeno con vocabulario técnico (\emph{``aquí falta validación de email''}).
\end{infoband}

\puente{Tienes el método. Toca aplicarlo: identificas un proceso real que se beneficiaría de un formulario (matrícula a una actividad, pedido a un negocio del barrio, inscripción a un evento, registro de asistencia), lo diseñas en Google Forms, lo pruebas tú mismo, lo enlazas a Sheets.}

\milcestacion{milcMagenta}{Estación 3 de 3 · Hacer}

\begin{softbox}{milcMagenta}{milcGris}{Cómo vas a construir tu producto}
\textbf{Actividad 3 · CREA --- Tu formulario funcional} (30 min · individual). Hora de crear. Elige un proceso real (matrícula a actividad, pedido a un negocio del barrio, inscripción a evento, registro de asistencia), abre Google Forms (gratuito, sin instalación), diseña el formulario con validaciones y lógica condicional, enláchalo a hoja de cálculo y compártelo.
\end{softbox}

\milcpaso{1}{milcMagenta}{Tu formulario se descompone en 5 piezas: (1) título + descripción; (2) 8-10 campos con tipos apropiados; (3) validaciones obligatorias en al menos 3 campos críticos; (4) 1 lógica condicional (\emph{``si responde X, salta a Y''}); (5) hoja de cálculo enlazada activa.}
\milcpaso{2}{milcMagenta}{Sigue el patrón \textbf{``probar como usuario antes de publicar''}: completa tu propio formulario 2 veces, una respondiendo \emph{``todo bien''} y otra forzando errores (email malo, número fuera de rango, campos vacíos). Si las validaciones funcionan, el formulario sirve; si no, ajusta.}
\milcpaso{3}{milcMagenta}{Cada campo debe justificar su existencia: \emph{¿qué decisión se toma con este dato?}. Si no hay decisión asociada, el campo se elimina. La economía del formulario es ética del diseño.}
\milcpaso{4}{milcMagenta}{Algoritmo: (1) abre forms.google.com; (2) crea formulario con título; (3) agrega los campos con sus tipos; (4) marca obligatorios y validaciones; (5) activa lógica condicional desde la pestaña respectiva; (6) enlaza a hoja de cálculo desde Respuestas; (7) prueba como usuario; (8) comparte URL.}

\begin{drawbox}{milcMagenta}{Checklist antes de publicar el formulario}
\checkbox Título y descripción claros. \\[1mm] \checkbox 8-10 campos máximo, cada uno justificado. \\[1mm] \checkbox Validaciones activas en email, número y campos obligatorios. \\[1mm] \checkbox Al menos 1 lógica condicional implementada. \\[1mm] \checkbox Hoja de cálculo enlazada y recibiendo respuestas. \\[1mm] \checkbox Probé el formulario como usuario y se completa en menos de 3 minutos. \\[1mm] \checkbox URL pública lista para compartir.
\end{drawbox}

\begin{softbox}{milcMagenta}{milcGris}{Plantilla de trabajo}
\textbf{Proceso que automatizo.} \textit{Nombre + responsable + dónde se usa.} \\[1mm] \textbf{Información que necesito.} \textit{Lista de 8-10 datos.} \\[1mm] \textbf{Campos del formulario.} \textit{Tipo de cada campo + validación + condicional si aplica.} \\[1mm] \textbf{Destino.} \textit{Hoja de cálculo enlazada + responsable de revisar respuestas.} \\[1mm] \textbf{Prueba personal.} \textit{Tiempo que tardé respondiéndolo + ajustes que hice tras la prueba.}
\end{softbox}

\begin{infoband}{milcMagenta}


\textbf{Lo que va al cuaderno (Actividad 3 · CREA).} Título: «Actividad 3 --- Mi formulario funcional». \textbf{Formato:} pantallazo del formulario + URL pública + pantallazo de la hoja con primeras respuestas + checklist marcado. \textbf{Extensión:} 1 página de cuaderno. \textbf{Sabes que terminaste cuando} compartiste el formulario con un compañero o familiar y recibiste al menos 1 respuesta real en la hoja.
\end{infoband}

\puente{Lo que sigue resume \emph{qué} entregas hoy: 1 formulario funcional + 1 hoja de cálculo conectada + URL pública. El criterio principal: que un usuario real pueda \textbf{responderlo en menos de 3 minutos} sin pedir ayuda.}

\milcestacion{milcMostaza}{Producto de la sesión}
\textbf{1 formulario digital funcional + hoja de cálculo enlazada}.
\begin{softbox}{milcMostaza}{milcGris}{Criterios de calidad}
(1) 8-10 campos con tipos apropiados. (2) Validaciones activas en email, número y obligatorios. (3) Al menos 1 lógica condicional implementada. (4) Hoja de cálculo enlazada activa. (5) URL pública compartible. (6) El formulario completo se responde en menos de 3 minutos.
\end{softbox}

\puente{Antes de cerrar, mira el formulario desde las cinco dimensiones humanas. Diseñar bien es respeto por el tiempo del usuario; diseñar mal es robo silencioso de los minutos de muchos.}

\milcestacion{milcVino}{Evaluación liberadora · cinco dimensiones}

\begin{softbox}{milcVino}{milcGris}{Sentido de la evaluación}
Evaluar no es solo revisar una nota. Es reconocer qué aprendí, qué cambió en mí, qué emociones manejé, cómo me relaciono con la ciudad y la comunidad, y qué le dejaré a quien viene detrás.
\end{softbox}

\textbf{Mi reflexión liberadora en cinco dimensiones.} Completa cada frase iniciada:

\small
\begin{tabularx}{\textwidth}{>{\bfseries\RaggedRight\arraybackslash}p{3.4cm}Y>{\RaggedRight\arraybackslash}p{4.6cm}}
\rowcolor{milcVino}\color{white}Dimensión & \color{white}\bfseries Mi respuesta (frame starter) & \color{white}\bfseries Algo que descubrí\\
1. Personal & Lo que más cambió en mí fue \linead{6.5cm} & \linead{4.2cm}\\
2. Emocional & La emoción más fuerte fue \linead{2.5cm} y la regulé \linead{3cm} & \linead{4.2cm}\\
3. Ciudadana & En democracia esto importa porque \linead{6cm} & \linead{4.2cm}\\
4. Local (Cartago) & En mi barrio sería útil para \linead{6.5cm} & \linead{4.2cm}\\
5. Intergeneracional & A mi abuela/abuelo le diría \linead{6.5cm} & \linead{4.2cm}\\
\end{tabularx}
\normalsize

\begin{infoband}{milcVino}
\textbf{Para la próxima sustentación.} Vuelve a esta tabla en seis meses. Verás que las respuestas cambian — no porque lo aprendido cambie, sino porque tú cambias. La evaluación liberadora es una conversación contigo a través del tiempo.
\end{infoband}


\puente{Y para terminar, tres voces sobre la captura de información. Dussel sobre los formularios que excluyen a quien no tiene celular, Séneca sobre la economía como virtud del diseño, Floridi sobre los formularios como infraestructura de la administración digital.}

\milcestacion{milcGradeDark}{Triángulo de pensamiento · cierre filosófico}

\begin{softbox}{milcVino}{milcGris}{Dussel · lente del nosotros}
\textit{Todo formulario decide quién puede responder y quién queda fuera por no tener el dispositivo o el conocimiento.}

Un formulario digital excluye automáticamente a quien no tiene celular con datos, a quien no maneja Google, a quien tiene baja alfabetización digital. La phronesis del diseño exige preguntar siempre: \emph{¿quién no podrá responder este formulario?}, y ofrecer un canal alternativo (papel, llamada, ayuda presencial) para incluir a quienes el medio digital deja fuera. Automatizar no es excusa para excluir.

\textbf{¿Quién no podría responder mi formulario digital? ¿Qué canal complementario habilitaría para que la información que pido no excluya a quienes más podrían necesitar el servicio?}
\end{softbox}
\linead{\linewidth}\\[1mm]
\linead{\linewidth}

\begin{softbox}{milcMostaza}{milcGris}{Estoicismo · Séneca · cuidado interior}
\textit{La economía de palabras es respeto por el tiempo de quien lee.}

Es tentador agregar campos \emph{``por si acaso''}. La sabiduría estoica recomienda lo contrario: \emph{cada pregunta que no usas te roba tiempo a ti como diseñador y a cada usuario que la responde}. Un formulario de 8 preguntas bien diseñadas es más útil que uno de 25 que nadie completa. La economía es virtud del oficio.

\textbf{¿Qué preguntas de mi formulario están \emph{``por si acaso''} y nunca se usarán? ¿Qué pasaría si las elimino?}
\end{softbox}
\linead{\linewidth}\\[1mm]
\linead{\linewidth}

\begin{softbox}{milcTurquesa}{milcGris}{Floridi · lente de la infoesfera}
\textit{Los formularios digitales son la infraestructura administrativa invisible de la sociedad contemporánea.}

Inscripciones a universidad, trámites estatales, compras en línea, atención en salud: casi toda interacción institucional contemporánea pasa por un formulario digital. Aprender a diseñarlos a los 16 años es alfabetización profesional anticipada: el lenguaje con el que las organizaciones reciben y procesan información de la ciudadanía.

\textbf{¿Cuántos formularios digitales respondí esta semana? ¿Cuántos estaban bien diseñados y cuántos me hicieron perder tiempo? ¿Qué aprendí sobre diseño respondiéndolos?}
\end{softbox}
\linead{\linewidth}\\[1mm]
\linead{\linewidth}

\begin{infoband}{milcNegro}
\textbf{Nota para el docente.} Las tres formulaciones del triángulo son la síntesis del autor de la guía, no citas textuales de los pensadores. Si vas a citarlos en clase, remite a la obra original.
\end{infoband}

\milcestacion{milcVino}{Mi autoevaluación · marca una casilla por fila}

{\small
\setlength{\extrarowheight}{2.2pt}
\begin{tabularx}{\textwidth}{>{\RaggedRight\arraybackslash\bfseries}p{3.0cm}YYY}
\rowcolor{milcVino}\color{white}\bfseries Criterio & \color{white}\bfseries Lo logré & \color{white}\bfseries En proceso & \color{white}\bfseries Necesito apoyo\\[.6mm]
Comprensión del tema & \milcopcion{Explico las ideas con mis palabras.} & \milcopcion{Reconozco las ideas, falta precisión.} & \milcopcion{Necesito repasar lo básico.}\\[.8mm]
\rowcolor{milcGris}Pensamiento computacional & \milcopcion{Usé los pilares al armar mi producto.} & \milcopcion{Usé algunos pilares.} & \milcopcion{Necesito releer la estación 2.}\\[.8mm]
Aplicación y producto & \milcopcion{Mi producto cumple los criterios.} & \milcopcion{Está armado, con ajustes pendientes.} & \milcopcion{Aún está en construcción.}\\[.8mm]
\rowcolor{milcGris}Reflexión 5 dimensiones & \milcopcion{Respondí cada una con honestidad.} & \milcopcion{Respondí algunas con detalle.} & \milcopcion{Mis respuestas son muy generales.}\\[.8mm]
Triángulo de pensamiento & \milcopcion{Conecté las 3 voces con mi trabajo.} & \milcopcion{Conecté al menos una.} & \milcopcion{No las relaciono con mi proceso.}\\[.8mm]
\end{tabularx}}

\vspace{3mm}
\textbf{Hoy entendí que:}\\[1mm]
\linead{\linewidth}\\[1mm]
\linead{\linewidth}

\textbf{Mi compromiso para la próxima sesión es:}\\[1mm]
\linead{\linewidth}\\[1mm]
\linead{\linewidth}

\begin{softbox}{milcMostaza}{milcGris}{Cita de despedida · Marco Aurelio}
\textit{``El obstáculo en el camino se vuelve el camino. Lo que estorba al camino se vuelve camino.''}\\[1mm]
Cada error de hoy, bien observado, se convierte en pieza del camino que pisarás mañana.
\end{softbox}

\small
\begin{tabularx}{\textwidth}{>{\bfseries}p{3.6cm}Y}
\rowcolor{milcGradeDark!12}Nombre completo: & \linead{10cm}\\
Fecha: & \linead{4cm} \quad Curso/grupo: \linead{4cm}\\
Mi firma: & \linead{9cm}\\
\end{tabularx}
\normalsize

\begin{infoband}{milcGradeDark}
\textbf{Para la próxima semana.} Pide a 3 personas reales (familiares, vecinos, compañeros) que respondan tu formulario, mide cuánto tarda cada uno y pregúntales qué les confundió. Ajusta el formulario con base en lo que digan. La phronesis del diseño se entrena escuchando a quien usa, no defendiendo lo que se hizo.
\end{infoband}

\end{document}
